
\includegraphics[width=8.28357in,height=1.96528in]{media/image1.jpeg}

\textbf{PHẦN I. Câu trắc nghiệm nhiều phương án lựa chọn.} Thí sinh trả
lời từ câu 1 đến câu 12. Mỗi câu hỏi thí sinh chỉ chọn một phương án.

\textbf{Câu 1. {[}HIMA{]}} Phương trình
\pandocbounded{\includegraphics[keepaspectratio]{media/image2.wmf}} với
\pandocbounded{\includegraphics[keepaspectratio]{media/image3.wmf}} có
nghiệm là:

\textbf{A.}
\pandocbounded{\includegraphics[keepaspectratio]{media/image4.wmf}} và
\pandocbounded{\includegraphics[keepaspectratio]{media/image5.wmf}}
\textbf{B.}
\pandocbounded{\includegraphics[keepaspectratio]{media/image6.wmf}} và
\pandocbounded{\includegraphics[keepaspectratio]{media/image7.wmf}}

\textbf{C.}
\pandocbounded{
\includegraphics[keepaspectratio]{media/image8.wmf}} và
\pandocbounded{\includegraphics[keepaspectratio]{media/image9.wmf}}
\textbf{D.}
\pandocbounded{\includegraphics[keepaspectratio]{media/image10.wmf}}
và\pandocbounded{\includegraphics[keepaspectratio]{media/image11.wmf}}

\textbf{\hl{Lời giải: B}}

\pandocbounded{\includegraphics[keepaspectratio]{media/image12.wmf}}

\[\Leftrightarrow \left\lbrack \begin{aligned}
 & x + 10^{o} = 30^{o} + k360^{o} \\
 & x + 10^{o} = 180^{o} - 30^{o} + k360^{o}
\end{aligned} \right.\  \Leftrightarrow \left\lbrack \begin{aligned}
 & x = 20^{o} + k360^{o} \\
 & x = 140^{o} + k360^{o}
\end{aligned} \right.\ \left( k \in \mathbb{Z} \right)\]

Mà \pandocbounded{\includegraphics[keepaspectratio]{media/image13.wmf}}.

\textbf{Câu 2. {[}HIMA{]}} Cho hàm số
\pandocbounded{\includegraphics[keepaspectratio]{media/image14.wmf}} xác
định trên \(\mathbb{R}\) và có đồ thị bên dưới.

\pandocbounded{\includegraphics[keepaspectratio]{media/image15.png}}

Mệnh đề nào dưới đây sai?

\textbf{A.} Hàm số đã cho có điểm cực tiểu bằng
\pandocbounded{\includegraphics[keepaspectratio]{media/image16.wmf}}.

\textbf{B.} Hàm số đã cho có điểm cực đại bằng
\pandocbounded{\includegraphics[keepaspectratio]{media/image17.wmf}}.

\textbf{C.} Hàm số đã cho có giá trị cực đại bằng
\pandocbounded{\includegraphics[keepaspectratio]{media/image18.wmf}}.

\textbf{D.} Hàm số đã cho có giá trị cực tiểu bằng
\pandocbounded{\includegraphics[keepaspectratio]{media/image19.wmf}}.

\textbf{\hl{Lời giải: C}}

Từ đồ thị đã cho, ta có

Hàm số đã cho có điểm cực tiểu bằng
\pandocbounded{\includegraphics[keepaspectratio]{media/image20.wmf}} và
giá trị cực tiểu bằng
\pandocbounded{\includegraphics[keepaspectratio]{media/image21.wmf}}.

Hàm số đã cho có điểm cực đại bằng
\pandocbounded{\includegraphics[keepaspectratio]{media/image22.wmf}} và
giá trị cực đại bằng
\pandocbounded{\includegraphics[keepaspectratio]{media/image23.wmf}}.

\textbf{Câu 3. {[}HIMA{]}} Cho cấp số cộng
\pandocbounded{\includegraphics[keepaspectratio]{media/image24.wmf}} với
\pandocbounded{\includegraphics[keepaspectratio]{media/image25.wmf}}.
Công sai của cấp số cộng đã cho bằng

\textbf{A.}
\pandocbounded{\includegraphics[keepaspectratio]{media/image26.wmf}}.
\textbf{B.}
\pandocbounded{\includegraphics[keepaspectratio]{media/image27.wmf}}.
\textbf{C.}
\pandocbounded{\includegraphics[keepaspectratio]{media/image28.wmf}}.
\textbf{D.}
\pandocbounded{\includegraphics[keepaspectratio]{media/image29.wmf}}.

\hl{\textbf{Lời} \textbf{giải: B}}

\textbf{Câu 4. {[}HIMA{]}} Hằng ngày ông Thắng đều đi xe buýt từ nhà đến
cơ quan. Dưới đây là bảng thống kê thời gian của 100 lần ông Thắng đi xe
buýt từ nhà đến cơ quan.

{\def\LTcaptype{none} % do not increment counter
\begin{longtable}[]{@{}
  >{\centering\arraybackslash}p{(\linewidth - 12\tabcolsep) * \real{0.1927}}
  >{\centering\arraybackslash}p{(\linewidth - 12\tabcolsep) * \real{0.1252}}
  >{\centering\arraybackslash}p{(\linewidth - 12\tabcolsep) * \real{0.1310}}
  >{\centering\arraybackslash}p{(\linewidth - 12\tabcolsep) * \real{0.1310}}
  >{\centering\arraybackslash}p{(\linewidth - 12\tabcolsep) * \real{0.1330}}
  >{\centering\arraybackslash}p{(\linewidth - 12\tabcolsep) * \real{0.1251}}
  >{\centering\arraybackslash}p{(\linewidth - 12\tabcolsep) * \real{0.1251}}@{}}
\toprule\noalign{}
\endhead
\bottomrule\noalign{}
\endlastfoot
Thời gian (phút) &
\pandocbounded{\includegraphics[keepaspectratio]{media/image30.wmf}} &
\pandocbounded{\includegraphics[keepaspectratio]{media/image31.wmf}} &
\pandocbounded{\includegraphics[keepaspectratio]{media/image32.wmf}} &
\pandocbounded{\includegraphics[keepaspectratio]{media/image33.wmf}} &
\pandocbounded{\includegraphics[keepaspectratio]{media/image34.wmf}} &
\pandocbounded{\includegraphics[keepaspectratio]{media/image35.wmf}} \\
Số lượt & 22 & 38 & 27 & 8 & 4 & 1 \\
\end{longtable}
}

Tứ phân vị thứ nhất của mẫu số liệu ghép nhóm trên (làm tròn kết quả đến
hàng phần trăm) là

\textbf{A.}
\pandocbounded{\includegraphics[keepaspectratio]{media/image36.wmf}}.
\textbf{B.}
\pandocbounded{\includegraphics[keepaspectratio]{media/image37.wmf}}.
\textbf{C.}
\pandocbounded{\includegraphics[keepaspectratio]{media/image38.wmf}}.
\textbf{D.}
\pandocbounded{\includegraphics[keepaspectratio]{media/image39.wmf}}.

\textbf{\hl{Lời giải: C}}

\textbf{Câu 5.} \textbf{{[}HIMA{]}} Cho hình lập phương
\pandocbounded{\includegraphics[keepaspectratio]{media/image40.wmf}}.

\includegraphics[width=2in,height=1.72708in,alt={Description: A black background with blue dots Description automatically generated}]{media/image41.png}

Góc giữa hai vectơ
\pandocbounded{\includegraphics[keepaspectratio]{media/image42.wmf}} và
\pandocbounded{\includegraphics[keepaspectratio]{media/image43.wmf}} là:

\textbf{A.}
\pandocbounded{\includegraphics[keepaspectratio]{media/image44.wmf}}.
\textbf{B.}
\pandocbounded{\includegraphics[keepaspectratio]{media/image45.wmf}}.
\textbf{C.}
\pandocbounded{\includegraphics[keepaspectratio]{media/image46.wmf}}.
\textbf{D.}
\pandocbounded{\includegraphics[keepaspectratio]{media/image47.wmf}}.

\textbf{\hl{Lời giải: C}}

\textbf{Câu 6.} \textbf{{[}HIMA{]}} Cho hàm số
\pandocbounded{\includegraphics[keepaspectratio]{media/image14.wmf}} xác
định trên \(\mathbb{R}\backslash\left\{ - 1 \right\}\) và bảng biến
thiên như sau:~

\includegraphics[width=3.81826in,height=1.3432in,alt={Description: A math equations and numbers Description automatically generated with medium confidence}]{media/image48.png}

Mệnh đề nào dưới đây đúng?

\textbf{A.} Hàm số đã cho nghịch biến trên các
khoảng\pandocbounded{\includegraphics[keepaspectratio]{media/image49.wmf}}
và \pandocbounded{\includegraphics[keepaspectratio]{media/image50.wmf}}.

\textbf{B.} Hàm số đã cho đồng biến trên khoảng
\pandocbounded{\includegraphics[keepaspectratio]{media/image51.wmf}}.

\textbf{C.} Hàm số đã cho có điểm cực tiểu bằng
\pandocbounded{\includegraphics[keepaspectratio]{media/image52.wmf}}.

\textbf{D.} Hàm số đã cho có giá trị cực đại bằng
\pandocbounded{\includegraphics[keepaspectratio]{media/image53.wmf}}.

\textbf{\hl{Lời giải: A}}

\textbf{Câu 7.} \textbf{{[}HIMA{]}} Cho hình chóp
\pandocbounded{\includegraphics[keepaspectratio]{media/image54.wmf}} có
đáy là hình thang,
\pandocbounded{\includegraphics[keepaspectratio]{media/image55.wmf}}//\pandocbounded{\includegraphics[keepaspectratio]{media/image56.wmf}}
và \pandocbounded{\includegraphics[keepaspectratio]{media/image57.wmf}}.
Lấy
\pandocbounded{\includegraphics[keepaspectratio]{media/image58.wmf}}thuộc
cạnh
\pandocbounded{\includegraphics[keepaspectratio]{media/image59.wmf}},
\pandocbounded{\includegraphics[keepaspectratio]{media/image60.wmf}}
thuộc cạnh
\pandocbounded{\includegraphics[keepaspectratio]{media/image61.wmf}} sao
cho
\pandocbounded{\includegraphics[keepaspectratio]{media/image62.wmf}}.
Khẳng định nào dưới đây đúng?

\textbf{A.} Đường thẳng
\pandocbounded{\includegraphics[keepaspectratio]{media/image63.wmf}}
song song với mặt phẳng
\pandocbounded{\includegraphics[keepaspectratio]{media/image64.wmf}}.

\textbf{B.} Đường thẳng
\pandocbounded{\includegraphics[keepaspectratio]{media/image65.wmf}} cắt
đường thẳng
\pandocbounded{\includegraphics[keepaspectratio]{media/image66.wmf}}.

\textbf{C.} Đường thẳng
\pandocbounded{\includegraphics[keepaspectratio]{media/image67.wmf}}
song song với mặt phẳng
\pandocbounded{\includegraphics[keepaspectratio]{media/image68.wmf}}.

\textbf{D.} Đường thẳng
\pandocbounded{\includegraphics[keepaspectratio]{media/image69.wmf}}
song song với mặt phẳng
\pandocbounded{\includegraphics[keepaspectratio]{media/image70.wmf}}.

\textbf{\hl{Lời giải: C}}

\includegraphics[width=2.59595in,height=1.98611in]{media/image71.png}

Vì \pandocbounded{\includegraphics[keepaspectratio]{media/image72.wmf}}
nên đường thẳng
\pandocbounded{\includegraphics[keepaspectratio]{media/image73.wmf}} //
\pandocbounded{\includegraphics[keepaspectratio]{media/image74.wmf}}. Mà
\pandocbounded{\includegraphics[keepaspectratio]{media/image75.wmf}},
\pandocbounded{\includegraphics[keepaspectratio]{media/image76.wmf}} nên
\pandocbounded{\includegraphics[keepaspectratio]{media/image77.wmf}}
song song với mặt phẳng
\pandocbounded{\includegraphics[keepaspectratio]{media/image78.wmf}}.

\textbf{Câu 8.} \textbf{{[}HIMA{]}} Nếu
\pandocbounded{\includegraphics[keepaspectratio]{media/image79.wmf}} thì
\pandocbounded{\includegraphics[keepaspectratio]{media/image80.wmf}}
bằng bao nhiêu?

\textbf{A.}
\pandocbounded{\includegraphics[keepaspectratio]{media/image81.wmf}}.
\textbf{B.}
\pandocbounded{\includegraphics[keepaspectratio]{media/image82.wmf}}.
\textbf{C.}
\pandocbounded{\includegraphics[keepaspectratio]{media/image83.wmf}}.
\textbf{D.}
\pandocbounded{\includegraphics[keepaspectratio]{media/image84.wmf}}.

\textbf{\hl{Lời giải: A}}

Ta có
\pandocbounded{\includegraphics[keepaspectratio]{media/image85.wmf}}.

\textbf{Câu 9.} \textbf{{[}HIMA{]}} Cho hình chóp
\pandocbounded{\includegraphics[keepaspectratio]{media/image86.wmf}}\textbf{,}
có đáy
\pandocbounded{\includegraphics[keepaspectratio]{media/image87.wmf}} là
hình bình hành tâm
\pandocbounded{\includegraphics[keepaspectratio]{media/image88.wmf}}.
Tam giác
\pandocbounded{\includegraphics[keepaspectratio]{media/image89.wmf}} đều
cạnh bằng
\pandocbounded{\includegraphics[keepaspectratio]{media/image90.wmf}}.

\includegraphics[width=2.6875in,height=1.84963in]{media/image91.png}

Khi đó
\pandocbounded{\includegraphics[keepaspectratio]{media/image92.wmf}}
bằng:

\textbf{A.}
\pandocbounded{\includegraphics[keepaspectratio]{media/image93.wmf}}.
\textbf{B.}
\pandocbounded{\includegraphics[keepaspectratio]{media/image94.wmf}}.
\textbf{C.}
\pandocbounded{\includegraphics[keepaspectratio]{media/image95.wmf}}.
\textbf{D.}
\pandocbounded{\includegraphics[keepaspectratio]{media/image96.wmf}}.

\textbf{\hl{Lời giải: D}}

Vì tam giác
\pandocbounded{\includegraphics[keepaspectratio]{media/image97.wmf}} đều
cạnh bằng
\pandocbounded{\includegraphics[keepaspectratio]{media/image98.wmf}} nên
\pandocbounded{\includegraphics[keepaspectratio]{media/image99.wmf}}.

Do đó
\pandocbounded{\includegraphics[keepaspectratio]{media/image100.wmf}}

\textbf{Câu 10.} \textbf{{[}HIMA{]}} Cho hình phẳng
\pandocbounded{\includegraphics[keepaspectratio]{media/image101.wmf}}
được giới hạn bởi đồ thị của hàm số
\pandocbounded{\includegraphics[keepaspectratio]{media/image102.wmf}} và
đồ thị của hàm số
\pandocbounded{\includegraphics[keepaspectratio]{media/image103.wmf}}
(tham khảo hình vẽ). Thể tích của khối tròn xoay thu được khi quay
\pandocbounded{\includegraphics[keepaspectratio]{media/image104.wmf}}
quanh trục
\pandocbounded{\includegraphics[keepaspectratio]{media/image105.wmf}}
bằng

\includegraphics[width=2.28472in,height=1.68805in,alt={Description: Ảnh có chứa biểu đồ Mô tả được tạo tự động}]{media/image106.png}

\textbf{A.}
\pandocbounded{\includegraphics[keepaspectratio]{media/image107.wmf}}.
\textbf{B.}
\pandocbounded{\includegraphics[keepaspectratio]{media/image108.wmf}}.
\textbf{C.}
\pandocbounded{\includegraphics[keepaspectratio]{media/image109.wmf}}.
\textbf{D.}
\pandocbounded{\includegraphics[keepaspectratio]{media/image110.wmf}}.

\hl{\textbf{Lời} \textbf{giải: B}}

\textbf{Câu 11.} \textbf{{[}HIMA{]}} Trong không gian tọa độ
\pandocbounded{\includegraphics[keepaspectratio]{media/image111.wmf}},
mặt phẳng
\pandocbounded{\includegraphics[keepaspectratio]{media/image112.wmf}}
vuông góc với mặt phẳng nào dưới đây?

\textbf{A.}
\pandocbounded{\includegraphics[keepaspectratio]{media/image113.wmf}}.
\textbf{B.}
\pandocbounded{\includegraphics[keepaspectratio]{media/image114.wmf}}.

\textbf{C.}
\pandocbounded{\includegraphics[keepaspectratio]{media/image115.wmf}}.
\textbf{D.}
\pandocbounded{\includegraphics[keepaspectratio]{media/image116.wmf}}.

\textbf{\hl{Lời giải: A}}

\textbf{Câu 12.} \textbf{{[}HIMA{]}} Trong không gian với hệ toạ độ
\pandocbounded{\includegraphics[keepaspectratio]{media/image117.wmf}},
khoảng cách từ điểm
\pandocbounded{\includegraphics[keepaspectratio]{media/image118.wmf}}
đến mặt phẳng
\pandocbounded{\includegraphics[keepaspectratio]{media/image119.wmf}}
bằng

\textbf{A.} 4. \textbf{B.} 5. \textbf{C.} 3. \textbf{D.} 2.

\textbf{\hl{Lời giải: D}}

\textbf{Phần II. Câu trắc nghiệm đúng sai.} Trong mỗi ý ở mỗi câu, hãy
chọn đúng hay sai

\textbf{Câu 1. {[}HIMA{]}} Bạn Nam chạy bộ trong 2 giờ, vận tốc
\pandocbounded{\includegraphics[keepaspectratio]{media/image120.wmf}}
phụ thuộc vào thời gian
\pandocbounded{\includegraphics[keepaspectratio]{media/image121.wmf}} có
đồ thị là một phần của đường Parabol với đỉnh
\pandocbounded{\includegraphics[keepaspectratio]{media/image122.wmf}} và
trục đối xứng song song với trục tung
\pandocbounded{\includegraphics[keepaspectratio]{media/image123.wmf}}
như hình vẽ.

\includegraphics[width=4.86662in,height=2.0587in]{media/image124.png}

\begin{enumerate}
\def\labelenumi{\alph{enumi})}
\item
  Vận tốc chạy bộ của bạn Nam trong giờ thứ
  \pandocbounded{\includegraphics[keepaspectratio]{media/image125.wmf}}tăng
  dần.
\item
  Vận tốc lớn nhất trong 2 giờ chạy của bạn Nam là
  \pandocbounded{\includegraphics[keepaspectratio]{media/image126.wmf}}.
\item
  Gia tốc chạy bộ của bạn Nam tại thời điểm
  \pandocbounded{\includegraphics[keepaspectratio]{media/image127.wmf}}giờ
  là
  \pandocbounded{\includegraphics[keepaspectratio]{media/image128.wmf}}.
\item
  Quãng đường bạn Nam chạy được trong 1 giờ 30 phút kể từ lúc bắt đầu
  chạy là
  \pandocbounded{\includegraphics[keepaspectratio]{media/image129.wmf}}
  (kết quả làm tròn đến 2 chữ số thập phân).
\end{enumerate}

\textbf{\hl{Lời giải: S-Đ-S-S}}

a) Từ đồ thị ta thấy: Vận tốc chạy bộ của bạn Nam trong khoảng
\pandocbounded{\includegraphics[keepaspectratio]{media/image130.wmf}}
giờ đầu tăng dần và
\pandocbounded{\includegraphics[keepaspectratio]{media/image130.wmf}}
giờ sau giảm dần.

Vậy vận tốc chạy bộ của bạn Nam trong trong giờ thứ
\pandocbounded{\includegraphics[keepaspectratio]{media/image125.wmf}}giảm
dần.

c) Đồ thị
\pandocbounded{\includegraphics[keepaspectratio]{media/image131.wmf}} đi
qua gốc tọa độ nên
\pandocbounded{\includegraphics[keepaspectratio]{media/image132.wmf}} có
dạng
\pandocbounded{\includegraphics[keepaspectratio]{media/image133.wmf}}.

Đồ thị
\pandocbounded{\includegraphics[keepaspectratio]{media/image134.wmf}} có
đỉnh là
\pandocbounded{\includegraphics[keepaspectratio]{media/image135.wmf}}
nên
\pandocbounded{\includegraphics[keepaspectratio]{media/image136.wmf}}

Gia tốc chạy bộ của bạn Nam tại thời điểm
\pandocbounded{\includegraphics[keepaspectratio]{media/image137.wmf}}
là~:
\pandocbounded{\includegraphics[keepaspectratio]{media/image138.wmf}}

\pandocbounded{\includegraphics[keepaspectratio]{media/image139.wmf}}

d) Ta có 1 giờ 30 phút = 1,5 giờ
\pandocbounded{\includegraphics[keepaspectratio]{media/image140.wmf}}.

Do đó, quãng đường bạn Nam chạy được trong 1 giờ 30 phút kể từ lúc bắt
đầu chạy là

\pandocbounded{\includegraphics[keepaspectratio]{media/image141.wmf}}.

\textbf{Câu 2. {[}HIMA{]}} Cho hàm số
\pandocbounded{\includegraphics[keepaspectratio]{media/image142.wmf}}

\begin{enumerate}
\def\labelenumi{\alph{enumi})}
\item
  Tập xác định của hàm số đã cho là
  \(D = \mathbb{R}\backslash\left\{ - \frac{1}{2} \right\}\).
\item
  Phương trình
  \pandocbounded{\includegraphics[keepaspectratio]{media/image143.wmf}}
  có hai nghiệm nguyên dương phân biệt.
\item
  Hàm số đã cho có hai điểm cực trị và hai điểm cực trị này nằm về hai
  phía của trục tung
\item
  Khi đồ thị hàm số đã cho có hai điểm cực trị thì đường thẳng đi qua
  hai điểm cực trị của đồ thị hàm số có phương trình
  \pandocbounded{\includegraphics[keepaspectratio]{media/image144.wmf}}
\end{enumerate}

\textbf{\hl{Lời giải: Đ-S-Đ-Đ}}

a) Tập xác định
\(D = \mathbb{R}\backslash\left\{ - \frac{1}{2} \right\}\).

b)
\pandocbounded{\includegraphics[keepaspectratio]{media/image145.wmf}},
\pandocbounded{\includegraphics[keepaspectratio]{media/image146.wmf}}\pandocbounded{\includegraphics[keepaspectratio]{media/image147.wmf}}.

c) Đồ thị hàm số có hai điểm cực trị là
\pandocbounded{\includegraphics[keepaspectratio]{media/image148.wmf}} và
\pandocbounded{\includegraphics[keepaspectratio]{media/image149.wmf}} và
hai điểm cực trị này nằm về hai phía của trục tung

d) Phương trình đường thẳng qua hai điểm cực trị
\pandocbounded{\includegraphics[keepaspectratio]{media/image150.wmf}}
của đồ thị hàm số đã cho là:
\pandocbounded{\includegraphics[keepaspectratio]{media/image151.wmf}}.

\textbf{Câu 3.} \textbf{{[}HIMA{]}} Thầy Nam muốn xây một ngôi nhà có
kích thước như hình 1. Ngôi nhà được gắn trong hệ trục tọa độ
\pandocbounded{\includegraphics[keepaspectratio]{media/image152.wmf}}
như hình 2 (đơn vị: mét), trong đó nền nhà, bốn bức tường là hình chữ
nhật và hai mái nhà là hai hình chữ nhật bằng nhau.

\includegraphics[width=5.39583in,height=2.25124in]{media/image153.png}

\begin{enumerate}
\def\labelenumi{\alph{enumi})}
\item
  \(\overrightarrow{OP} = (4;0;4,7)\)
\item
  \pandocbounded{\includegraphics[keepaspectratio]{media/image154.wmf}}
\item
  Thầy Nam muốn treo bóng đèn tại trung điểm cạnh
  \pandocbounded{\includegraphics[keepaspectratio]{media/image155.wmf}}.
  Tọa độ vị trí bóng đèn là
  \pandocbounded{\includegraphics[keepaspectratio]{media/image156.wmf}}.
\item
  Hai mái nhà hợp nhau một góc bằng
  \pandocbounded{\includegraphics[keepaspectratio]{media/image157.wmf}}
  (làm tròn kết quả đến hàng phần mười của độ).
\end{enumerate}

\textbf{\hl{Lời giải: Đ-Đ-Đ-S}}

Từ hình vẽ, ta có kích thước ngôi nhà như sau:

\pandocbounded{\includegraphics[keepaspectratio]{media/image158.wmf}}

\pandocbounded{\includegraphics[keepaspectratio]{media/image159.wmf}}

Tọa độ các điểm là :

\pandocbounded{\includegraphics[keepaspectratio]{media/image160.wmf}}.

\pandocbounded{\includegraphics[keepaspectratio]{media/image161.wmf}}

a) \pandocbounded{\includegraphics[keepaspectratio]{media/image162.wmf}}

b) \pandocbounded{\includegraphics[keepaspectratio]{media/image163.wmf}}

c) Tọa độ trung điểm cạnh
\pandocbounded{\includegraphics[keepaspectratio]{media/image155.wmf}} là
\pandocbounded{\includegraphics[keepaspectratio]{media/image156.wmf}}

d) góc hợp bởi hai mái nhà là góc hai mặt phẳng
\pandocbounded{\includegraphics[keepaspectratio]{media/image164.wmf}} và
( \pandocbounded{\includegraphics[keepaspectratio]{media/image165.wmf}}
hay góc
\pandocbounded{\includegraphics[keepaspectratio]{media/image166.wmf}}

Ta có:
\pandocbounded{\includegraphics[keepaspectratio]{media/image167.wmf}}.

Suy ra
\(\cos\widehat{EPF} = \cos(\overrightarrow{PE},\overrightarrow{PF}) = \frac{\overrightarrow{PE} \cdot \overrightarrow{PF}}{|\overrightarrow{PE}| \cdot |\overrightarrow{PF}|} = \frac{( - 4).(4) + 0.0 + ( - 1,2).( - 1,2)}{\sqrt{( - 4)^{2} + 0^{2} + ( - 1,2)^{2}} \cdot \sqrt{4^{2} + 0^{2} + ( - 1,2)^{2}}} = - \frac{91}{109}\text{. }\)

\(\Rightarrow \widehat{EPF} \approx 146,6^{0}\). Vậy hai mái nhà hợp
nhau một góc bằng
\pandocbounded{\includegraphics[keepaspectratio]{media/image168.wmf}}.

\textbf{Câu 4. {[}HIMA{]}} Trong một cuộc khảo sát tình trạng công việc
trên
\pandocbounded{\includegraphics[keepaspectratio]{media/image169.wmf}}
người chỉ có bằng tốt nghiệp THPT tại một địa phương, người ta thu được
số liệu như bảng dưới đây

\includegraphics[width=2.96255in,height=0.86767in]{media/image170.png}

Chọn ngẫu nhiên một người trong nhóm này. Khi đó

\begin{enumerate}
\def\labelenumi{\alph{enumi})}
\item
  Xác suất để chọn được một nam là
  \pandocbounded{\includegraphics[keepaspectratio]{media/image171.wmf}}.
\item
  Xác suất để chọn được một người có việc làm là
  \pandocbounded{\includegraphics[keepaspectratio]{media/image172.wmf}}.
\item
  Tại địa phương này, nếu chỉ có bằng tốt nghiệp THPT thì tỉ lệ nữ thất
  nghiệp sẽ cao hơn nam. Khảo sát cho thấy xác suất để một người thất
  nghiệp khi người đó là nữ cao gấp 7 lần xác suất để một người thất
  nghiệp khi người đó là nam.
\item
  Biết rằng đã chọn được một người có việc làm, xác suất để người này là
  nữ là
  \pandocbounded{\includegraphics[keepaspectratio]{media/image173.wmf}}.
\end{enumerate}

\textbf{\hl{Lời giải: Đ-Đ-S-Đ}}

a) Gọi
\pandocbounded{\includegraphics[keepaspectratio]{media/image174.wmf}} là
biến cố:''Chọn được một người là nam''.

Ta có
\pandocbounded{\includegraphics[keepaspectratio]{media/image175.wmf}}

\pandocbounded{\includegraphics[keepaspectratio]{media/image176.wmf}}

\pandocbounded{\includegraphics[keepaspectratio]{media/image177.wmf}}

b) Gọi
\pandocbounded{\includegraphics[keepaspectratio]{media/image178.wmf}} là
biến cố:''Chọn được một người có việc làm''.

Ta có
\pandocbounded{\includegraphics[keepaspectratio]{media/image179.wmf}}

\pandocbounded{\includegraphics[keepaspectratio]{media/image180.wmf}}

\pandocbounded{\includegraphics[keepaspectratio]{media/image181.wmf}}

c) Xác suất để một người thất nghiệp khi người đó là nữ cao gấp
\pandocbounded{\includegraphics[keepaspectratio]{media/image182.wmf}}
lần xác suất để một người thất nghiệp khi người đó là nam.

d) Gọi
\pandocbounded{\includegraphics[keepaspectratio]{media/image183.wmf}} là
biến cố:''Chọn được một người là nữ''.

Gọi
\pandocbounded{\includegraphics[keepaspectratio]{media/image184.wmf}} là
biến cố:''Chọn được một người có việc làm''.

Ta tính
\pandocbounded{\includegraphics[keepaspectratio]{media/image185.wmf}}

Ta có
\pandocbounded{\includegraphics[keepaspectratio]{media/image186.wmf}}

\textbf{PHẦN III. Câu trắc nghiệm trả lời ngắn.} Thí sinh trả lời từ câu
1 đến câu 6.

\textbf{Câu 1. {[}HIMA{]}} Để trang trí một bảng gỗ hình chữ nhật
\pandocbounded{\includegraphics[keepaspectratio]{media/image187.wmf}} có
chiều dài
\pandocbounded{\includegraphics[keepaspectratio]{media/image188.wmf}} và
chiều rộng
\pandocbounded{\includegraphics[keepaspectratio]{media/image189.wmf}},
người ta thiết kế một logo là hình phẳng giới hạn bởi nửa đường tròn
đường kính
\pandocbounded{\includegraphics[keepaspectratio]{media/image190.wmf}}tiếp
xúc với
\pandocbounded{\includegraphics[keepaspectratio]{media/image191.wmf}},
hai đường cong
\pandocbounded{\includegraphics[keepaspectratio]{media/image192.wmf}} là
một phần của các đường parabol có trục đối xứng lần lượt là
\pandocbounded{\includegraphics[keepaspectratio]{media/image193.wmf}}
với
\pandocbounded{\includegraphics[keepaspectratio]{media/image194.wmf}} là
trung điểm của
\pandocbounded{\includegraphics[keepaspectratio]{media/image195.wmf}}(tham
khảo hình vẽ).

\includegraphics[width=2.36111in,height=2.96369in]{media/image196.png}

Phần logo được sơn màu xanh với chi phí
\pandocbounded{\includegraphics[keepaspectratio]{media/image197.wmf}}
đồng/
\pandocbounded{\includegraphics[keepaspectratio]{media/image198.wmf}} và
phần còn lại của bảng gỗ được sơn màu trắng với chi phí
\pandocbounded{\includegraphics[keepaspectratio]{media/image199.wmf}}
đồng/
\pandocbounded{\includegraphics[keepaspectratio]{media/image200.wmf}}.
Hỏi cần bỏ ra bao nhiêu nghìn đồng để trang trí bảng gỗ trên (kết quả
làm tròn đến hàng đơn vị)?

\hl{\textbf{Lời} \textbf{giải: 1923}}

Diện tích hình chữ nhật
\pandocbounded{\includegraphics[keepaspectratio]{media/image201.wmf}}
bằng
\pandocbounded{\includegraphics[keepaspectratio]{media/image202.wmf}}.

Diện tích của nửa đường tròn đường kính
\pandocbounded{\includegraphics[keepaspectratio]{media/image203.wmf}}
bằng
\pandocbounded{\includegraphics[keepaspectratio]{media/image204.wmf}}.

Ta xác định phương trình của parabol.

Chọn hệ trục tọa độ
\pandocbounded{\includegraphics[keepaspectratio]{media/image205.wmf}}
như hình vẽ:

Gọi phương trình parabol có đỉnh
\pandocbounded{\includegraphics[keepaspectratio]{media/image206.wmf}} là
\pandocbounded{\includegraphics[keepaspectratio]{media/image207.wmf}},
theo bài ra ta có:

\pandocbounded{\includegraphics[keepaspectratio]{media/image208.wmf}}

Vậy
\pandocbounded{\includegraphics[keepaspectratio]{media/image209.wmf}}.

Diện tích của hình phẳng giới hạn bởi
\pandocbounded{\includegraphics[keepaspectratio]{media/image210.wmf}},
trục hoành và hai đường thẳng
\pandocbounded{\includegraphics[keepaspectratio]{media/image211.wmf}}
bằng:
\pandocbounded{\includegraphics[keepaspectratio]{media/image212.wmf}}.

Do tính đối xứng nên hình phẳng giới hạn bởi parabol có đỉnh
\pandocbounded{\includegraphics[keepaspectratio]{media/image213.wmf}},
trục đối xứng
\pandocbounded{\includegraphics[keepaspectratio]{media/image214.wmf}},
trục hoành và hai đường thẳng
\pandocbounded{\includegraphics[keepaspectratio]{media/image215.wmf}} có
diện tích bằng
\pandocbounded{\includegraphics[keepaspectratio]{media/image216.wmf}}.

Tổng diện tích phần bảng gỗ sơn màu xanh là:
\pandocbounded{\includegraphics[keepaspectratio]{media/image217.wmf}}.

Diện tích của phần bảng gỗ sơn màu trắng là:
\pandocbounded{\includegraphics[keepaspectratio]{media/image218.wmf}}.

Chi phí cần bỏ ra là:
\pandocbounded{\includegraphics[keepaspectratio]{media/image219.wmf}}
(nghìn đồng).

\textbf{Câu 2. {[}HIMA{]}} Một công ty chuyên sản xuất dụng cụ thể thao
nhận được đơn đặt hàng sản xuất
\pandocbounded{\includegraphics[keepaspectratio]{media/image220.wmf}}
quả bóng rổ. Công ty có một số máy móc, mỗi máy có khả năng sản xuất
\pandocbounded{\includegraphics[keepaspectratio]{media/image221.wmf}}
quả bóng rổ trong một giờ. Chi phí thiết lập mỗi máy là
\pandocbounded{\includegraphics[keepaspectratio]{media/image222.wmf}}
nghìn đồng. Sau khi thiết lập, quá trình sản xuất sẽ diễn ra hoàn toàn
tự động và chỉ cần có người giám sát. Chi phí trả cho người giám sát là
\pandocbounded{\includegraphics[keepaspectratio]{media/image223.wmf}}
nghìn đồng mỗi giờ. Số máy móc công ty cần sử dụng để chi phí hoạt động
đạt mức thấp nhất là bao nhiêu?

\hl{\textbf{Lời} \textbf{giải: 15}}

Gọi số máy móc công ty cần sử dụng để sản xuất
\pandocbounded{\includegraphics[keepaspectratio]{media/image224.wmf}}
quả bóng rổ là:
\pandocbounded{\includegraphics[keepaspectratio]{media/image225.wmf}}
(\pandocbounded{\includegraphics[keepaspectratio]{media/image226.wmf}}).

Chi phí thiết lập mỗi máy là
\pandocbounded{\includegraphics[keepaspectratio]{media/image227.wmf}}
nghìn đồng nên chi phí thiết lập
\pandocbounded{\includegraphics[keepaspectratio]{media/image228.wmf}}
máy là
\pandocbounded{\includegraphics[keepaspectratio]{media/image229.wmf}}nghìn
đồng.

Do mỗi máy sản xuất trong 1h được
\pandocbounded{\includegraphics[keepaspectratio]{media/image230.wmf}}
quả bóng rổ nên
\pandocbounded{\includegraphics[keepaspectratio]{media/image231.wmf}}
máy trong 1h sẽ sản xuất được
\pandocbounded{\includegraphics[keepaspectratio]{media/image232.wmf}}
quả.

Để sản xuất
\pandocbounded{\includegraphics[keepaspectratio]{media/image233.wmf}}
quả bóng thì cần
\pandocbounded{\includegraphics[keepaspectratio]{media/image234.wmf}}(h).

Chi phí trả cho người giám sát là
\pandocbounded{\includegraphics[keepaspectratio]{media/image235.wmf}}
nghìn đồng 1h nên chi phí cần trả cho người giám sát là:
\pandocbounded{\includegraphics[keepaspectratio]{media/image236.wmf}}nghìn
đồng.

Tổng chi phí phải trả đề sản xuất
\pandocbounded{\includegraphics[keepaspectratio]{media/image237.wmf}}
quả bóng là:
\pandocbounded{\includegraphics[keepaspectratio]{media/image238.wmf}}

Ta có:
\pandocbounded{\includegraphics[keepaspectratio]{media/image239.wmf}}

BBT:

\includegraphics[width=3.25in,height=1.16667in]{media/image240.png}

Từ BBT ta có
\pandocbounded{\includegraphics[keepaspectratio]{media/image241.wmf}}đạt
giá trị bé nhất tại
\pandocbounded{\includegraphics[keepaspectratio]{media/image242.wmf}}.

Vậy số máy móc công ty cần sử dụng để chi phí hoạt động đạt mức thấp
nhất là 15.

\textbf{Câu 3. {[}HIMA{]}} Cho hình lăng trụ đứng
\pandocbounded{\includegraphics[keepaspectratio]{media/image243.wmf}} có
đáy
\pandocbounded{\includegraphics[keepaspectratio]{media/image244.wmf}} là
tam giác vuông tại
\pandocbounded{\includegraphics[keepaspectratio]{media/image245.wmf}},
\pandocbounded{\includegraphics[keepaspectratio]{media/image246.wmf}} và
\pandocbounded{\includegraphics[keepaspectratio]{media/image247.wmf}}.
Tính khoảng cách từ điểm
\pandocbounded{\includegraphics[keepaspectratio]{media/image248.wmf}}
đến mặt phẳng
\pandocbounded{\includegraphics[keepaspectratio]{media/image249.wmf}}.

\hl{\textbf{Lời} \textbf{giải: 3}}

\includegraphics[width=2.18472in,height=1.94931in]{media/image250.emf}

Kẻ\pandocbounded{\includegraphics[keepaspectratio]{media/image251.wmf}};
\pandocbounded{\includegraphics[keepaspectratio]{media/image252.wmf}}.

Ta có:
\pandocbounded{\includegraphics[keepaspectratio]{media/image253.wmf}}.

Mà
\pandocbounded{\includegraphics[keepaspectratio]{media/image254.wmf}}.

Do đó khoảng cách từ điểm
\pandocbounded{\includegraphics[keepaspectratio]{media/image255.wmf}}
đến mặt phẳng
\pandocbounded{\includegraphics[keepaspectratio]{media/image256.wmf}}
bằng
\pandocbounded{\includegraphics[keepaspectratio]{media/image257.wmf}}.

Xét tam giác
\pandocbounded{\includegraphics[keepaspectratio]{media/image258.wmf}}
vuông tại
\pandocbounded{\includegraphics[keepaspectratio]{media/image259.wmf}} ta
có
\pandocbounded{\includegraphics[keepaspectratio]{media/image260.wmf}}.

Xét tam giác
\pandocbounded{\includegraphics[keepaspectratio]{media/image261.wmf}}
vuông tại
\pandocbounded{\includegraphics[keepaspectratio]{media/image262.wmf}} ta
có
\pandocbounded{\includegraphics[keepaspectratio]{media/image263.wmf}}.

\textbf{Câu 4. {[}HIMA{]}} Có hai hộp đựng các viên bi: Hộp
\pandocbounded{\includegraphics[keepaspectratio]{media/image264.wmf}} có
6 bi đỏ và 4 bi xanh, hộp
\pandocbounded{\includegraphics[keepaspectratio]{media/image265.wmf}} có
5 bi đỏ và 5 bi xanh (các viên bi có cùng kích thước và khối lượng). Bạn
Minh lấy ngẫu nhiên một viên bi từ hộp
\pandocbounded{\includegraphics[keepaspectratio]{media/image266.wmf}},
bạn Như lấy ngẫu nhiên một viên bi từ hộp
\pandocbounded{\includegraphics[keepaspectratio]{media/image267.wmf}}.
Biết rằng có viên bi đỏ được lấy ra, khi đó xác suất để bạn Như lấy được
viên bi đỏ là
\pandocbounded{\includegraphics[keepaspectratio]{media/image268.wmf}}(\pandocbounded{\includegraphics[keepaspectratio]{media/image269.wmf}}
tối giản). Tính
\pandocbounded{\includegraphics[keepaspectratio]{media/image270.wmf}}

\hl{\textbf{Lời} \textbf{giải: -3}}

Gọi
\pandocbounded{\includegraphics[keepaspectratio]{media/image271.wmf}} là
biến cố: ``Bạn Như lấy được viên bi đỏ''.

Gọi
\pandocbounded{\includegraphics[keepaspectratio]{media/image272.wmf}} là
biến cố: ``Có viên bi đỏ trong hai viên được lấy ra''.

Ta cần tìm xác suất
\pandocbounded{\includegraphics[keepaspectratio]{media/image273.wmf}}.

Theo công thức xác suất có điều kiện:
\pandocbounded{\includegraphics[keepaspectratio]{media/image274.wmf}}

Ta có
\pandocbounded{\includegraphics[keepaspectratio]{media/image275.wmf}}
(số cách lấy hai viên bi trong đó có ít nhất 1 bi đỏ) và
\pandocbounded{\includegraphics[keepaspectratio]{media/image276.wmf}}
(số cách lấy hai viên bi trong đó viên bi của bạn Như là bi đỏ).

Thay số, ta được
\pandocbounded{\includegraphics[keepaspectratio]{media/image277.wmf}}

Vậy
\pandocbounded{\includegraphics[keepaspectratio]{media/image278.wmf}}

\textbf{Câu 5. {[}HIMA{]}} Trong không gian với hệ trục
\pandocbounded{\includegraphics[keepaspectratio]{media/image279.wmf}},
giả sử mặt đất trùng với mặt phẳng
\pandocbounded{\includegraphics[keepaspectratio]{media/image280.wmf}}.
Một bóng đèn trang trí dạng khối cầu có tâm
\pandocbounded{\includegraphics[keepaspectratio]{media/image281.wmf}} và
bán kính
\pandocbounded{\includegraphics[keepaspectratio]{media/image282.wmf}}
được treo cố định lên trần nhà (đơn vị trên mỗi trục tọa độ là
\pandocbounded{\includegraphics[keepaspectratio]{media/image283.wmf}}).
Một con kiến bò tùy ý trên bóng đèn và một con kiến khác bò tùy ý trên
mặt đất, giả sử vectơ tạo bởi tọa độ vị trí của 2 con kiến luôn cùng
phương với đường thẳng
\pandocbounded{\includegraphics[keepaspectratio]{media/image284.wmf}}(coi
mỗi con kiến là một điểm). Biết lúc 2 con kiến gần nhau nhất có khoảng
cách bằng
\pandocbounded{\includegraphics[keepaspectratio]{media/image285.wmf}}.
Bán khối cầu có độ dài bao nhiêu
\pandocbounded{\includegraphics[keepaspectratio]{media/image286.wmf}}.

\includegraphics[width=1.95139in,height=2.08331in]{media/image287.png}

\hl{\textbf{Lời} \textbf{giải: 20}}

\includegraphics[width=2.40972in,height=1.91667in]{media/image288.png}

+ Gọi
\pandocbounded{\includegraphics[keepaspectratio]{media/image289.wmf}} là
góc tạo bởi
\pandocbounded{\includegraphics[keepaspectratio]{media/image290.wmf}} và
mặt
\pandocbounded{\includegraphics[keepaspectratio]{media/image291.wmf}}
suy ra:
\pandocbounded{\includegraphics[keepaspectratio]{media/image292.wmf}}.

Gọi
\pandocbounded{\includegraphics[keepaspectratio]{media/image293.wmf}} là
hình chiếu của
\pandocbounded{\includegraphics[keepaspectratio]{media/image294.wmf}}
lên
\pandocbounded{\includegraphics[keepaspectratio]{media/image295.wmf}},
gọi
\pandocbounded{\includegraphics[keepaspectratio]{media/image296.wmf}} là
giao điểm của
\pandocbounded{\includegraphics[keepaspectratio]{media/image297.wmf}} và
khối cầu
\pandocbounded{\includegraphics[keepaspectratio]{media/image298.wmf}}.
\pandocbounded{\includegraphics[keepaspectratio]{media/image299.wmf}} là
giao điểm của đường thẳng qua
\pandocbounded{\includegraphics[keepaspectratio]{media/image300.wmf}} và
song song
\pandocbounded{\includegraphics[keepaspectratio]{media/image301.wmf}}
với
\pandocbounded{\includegraphics[keepaspectratio]{media/image302.wmf}}.

Vì \pandocbounded{\includegraphics[keepaspectratio]{media/image303.wmf}}
là độ dài ngắn nhất giữa khối cầu
\pandocbounded{\includegraphics[keepaspectratio]{media/image304.wmf}} và
\pandocbounded{\includegraphics[keepaspectratio]{media/image305.wmf}}
nên lúc 2 con kiến gần nhau nhất thì một con ở
\pandocbounded{\includegraphics[keepaspectratio]{media/image306.wmf}} và
một con ở
\pandocbounded{\includegraphics[keepaspectratio]{media/image307.wmf}}.

+ Theo giả thiết:
\pandocbounded{\includegraphics[keepaspectratio]{media/image308.wmf}},
suy ra:
\pandocbounded{\includegraphics[keepaspectratio]{media/image309.wmf}}.

Vậy
\pandocbounded{\includegraphics[keepaspectratio]{media/image310.wmf}}

\textbf{Câu 6. {[}HIMA{]}} Một mô hình trang trí có dạng là hình lập
phương
\pandocbounded{\includegraphics[keepaspectratio]{media/image311.wmf}},
cạnh bằng
\pandocbounded{\includegraphics[keepaspectratio]{media/image312.wmf}}m
(\emph{như hình vẽ}). Người ta cần nối một đường dây điện đi từ điểm
\pandocbounded{\includegraphics[keepaspectratio]{media/image313.wmf}}
(là trung điểm của
\pandocbounded{\includegraphics[keepaspectratio]{media/image314.wmf}})
đi qua điểm
\pandocbounded{\includegraphics[keepaspectratio]{media/image315.wmf}}
thuộc cạnh
\pandocbounded{\includegraphics[keepaspectratio]{media/image316.wmf}},
đi tiếp qua điểm
\pandocbounded{\includegraphics[keepaspectratio]{media/image317.wmf}}
thuộc cạnh
\pandocbounded{\includegraphics[keepaspectratio]{media/image318.wmf}}
rồi tới điểm
\pandocbounded{\includegraphics[keepaspectratio]{media/image319.wmf}}.
Biết độ dài đoạn dây điện bằng
\pandocbounded{\includegraphics[keepaspectratio]{media/image320.wmf}}m.
Tính độ dài đoạn
\pandocbounded{\includegraphics[keepaspectratio]{media/image321.wmf}}(\emph{làm
tròn kết quả đến hàng phần trăm})

\includegraphics[width=3.54861in,height=2.52778in,alt={Description: A drawing of a cube with lines and letters AI-generated content may be incorrect.}]{media/image322.png}

\hl{\textbf{Lời} \textbf{giải: 4,17}}

Trải hình lập phương ra trên mặt phẳng ta được hình vẽ

\includegraphics[width=2.11806in,height=2.08522in,alt={Description: A diagram of a rectangle with letters and numbers Description automatically generated}]{media/image323.png}

Ta có
\pandocbounded{\includegraphics[keepaspectratio]{media/image324.wmf}} và
\pandocbounded{\includegraphics[keepaspectratio]{media/image325.wmf}}.

Ta thấy
\pandocbounded{\includegraphics[keepaspectratio]{media/image326.wmf}}

Suy ra
\pandocbounded{\includegraphics[keepaspectratio]{media/image327.wmf}}
thẳng hàng

Xét tam giác
\pandocbounded{\includegraphics[keepaspectratio]{media/image328.wmf}},
có

+
\(\tan\widehat{D'B'E} = \frac{15}{20} = \frac{A'N}{10}\)\pandocbounded{\includegraphics[keepaspectratio]{media/image329.wmf}}
và
\pandocbounded{\includegraphics[keepaspectratio]{media/image330.wmf}}.

+ \pandocbounded{\includegraphics[keepaspectratio]{media/image331.wmf}}
và
\pandocbounded{\includegraphics[keepaspectratio]{media/image332.wmf}}.

Vậy
\pandocbounded{\includegraphics[keepaspectratio]{media/image333.wmf}}(m)

\begin{quote}
\emph{Chị Hiền Mai chúc các em làm bài thật tốt!!}
\end{quote}

\textbf{-\/-\/-\/-\/-\/-\/-\/-\/-\/-\/-\/-\/-\/-\/-\/-
HẾT-\/-\/-\/-\/-\/-\/-\/-\/-\/-\/-\/-\/-\/-\/-\/-}
